\documentclass[14pt,a4paper]{extarticle}

% -------------------- Пакеты --------------------
\usepackage[utf8]{inputenc}
\usepackage[T2A]{fontenc}
\usepackage[russian]{babel}
\usepackage{geometry}
\geometry{
	left=25mm,
	right=20mm,
	top=20mm,
	bottom=20mm
}
\usepackage{setspace}
\onehalfspacing
\usepackage{indentfirst}
\usepackage{amsmath}
\usepackage{graphicx}
\usepackage{listings}
\usepackage{xcolor}

% Настройка оформления кода
\definecolor{codegray}{rgb}{0.5,0.5,0.5}
\definecolor{codebg}{rgb}{0.97,0.97,0.97}
\definecolor{keywordcolor}{rgb}{0.0,0.2,0.7}

\lstset{
	backgroundcolor=\color{codebg},
	basicstyle=\ttfamily\footnotesize,
	breaklines=true,
	captionpos=b,
	commentstyle=\color{codegray},
	keywordstyle=\color{keywordcolor}\bfseries,
	numbers=left,
	numbersep=6pt,
	numberstyle=\tiny\color{codegray},
	frame=single,
	rulecolor=\color{black!20},
	showstringspaces=false,
	tabsize=4
}

\begin{document}

% -------------------- Шапка --------------------
\begin{flushright}
Дисциплина: Компьютерная графика\\
Студент: Онищук Тимофей Дмитриевич\\
Группа: 3 курс, 2 группа\\
Вариант: Чётный (CMYK -- RGB -- HSV)
\end{flushright}

\vspace{1.0cm}

% -------------------- Заголовки --------------------
\begin{center}
\textbf{Отчёт по лабораторной работе №1}\\
\textbf{Тема: Цветовые модели: CMYK, RGB, HSV и взаимные преобразования}
\end{center}

\vspace{0.5cm}

% -------------------- Текст --------------------
\noindent\textbf{1. Цель работы}

Изучение физических принципов и математических моделей представления цвета (RGB, CMYK, HSV), реализация взаимных переходов между ними и разработка интерактивного графического приложения с мгновенным синхронным пересчётом параметров во всех трёх пространствах.

\vspace{0.5cm}
\noindent\textbf{2. Теоретические сведения и формулы преобразования}

\textbf{2.1. Модель RGB $\leftrightarrow$ CMYK}

Аддитивная модель RGB описывает излучаемый свет ($R, G, B \in [0, 255]$), а субтрактивная CMYK — отражённый ($C, M, Y, K \in [0, 1]$).
Пусть нормированные компоненты $R' = R / 255$, $G' = G / 255$, $B' = B / 255$. Тогда прямой перевод (RGB $\to$ CMYK) выполняется по формулам:
\begin{align*}
K &= 1 - \max(R', G', B'), \\
C &= \frac{1 - R' - K}{1 - K}, \quad M = \frac{1 - G' - K}{1 - K}, \quad Y = \frac{1 - B' - K}{1 - K} \quad (\text{при } K < 1).
\end{align*}
Если $R = G = B = 0$, то $C = M = Y = 0$, $K = 1$. Обратное преобразование (CMYK $\to$ RGB):
\[
R = 255 \cdot (1 - C) \cdot (1 - K), \quad G = 255 \cdot (1 - M) \cdot (1 - K), \quad B = 255 \cdot (1 - Y) \cdot (1 - K).
\]

\textbf{2.2. Модель RGB $\leftrightarrow$ HSV}

Модель HSV описывает цвет в цилиндрических координатах: цветовой тон $H \in [0^\circ, 360^\circ]$, насыщенность $S \in [0, 1]$ и яркость $V \in [0, 1]$.
Пусть $V = \max(R', G', B')$, $\Delta = V - \min(R', G', B')$.
\[
S = \begin{cases} 0, & V = 0 \\ \frac{\Delta}{V}, & V \neq 0 \end{cases}, \qquad
H = \begin{cases} 
0^\circ, & \Delta = 0 \\
60^\circ \cdot \left(\frac{G' - B'}{\Delta} \bmod 6\right), & V = R' \\
60^\circ \cdot \left(\frac{B' - R'}{\Delta} + 2\right), & V = G' \\
60^\circ \cdot \left(\frac{R' - G'}{\Delta} + 4\right), & V = B'
\end{cases}
\]

Обратный перевод из HSV в RGB выполняется вычислением вспомогательных величин $C = V \cdot S$, $X = C \cdot (1 - |(H / 60^\circ) \bmod 2 - 1|)$ и $m = V - C$, после чего промежуточные значения $(R_1, G_1, B_1)$ в зависимости от сектора угла $H$ смещаются на $m$ и масштабируются на 255.

\vspace{0.5cm}
\noindent\textbf{3. Описание интерфейса и способов задания цвета}

Приложение реализовано на языке Python с использованием библиотеки графического интерфейса \texttt{tkinter}. В соответствии с заданием обеспечены 3 способа ввода цвета:
\begin{enumerate}
    \item \textbf{Плавное изменение ползунками (Scale)}: для каждого из 10 каналов предусмотрены независимые горизонтальные шкалы.
    \item \textbf{Точный численный ввод (Entry)}: текстовые поля с валидацией ввода для прямого указания значений (RGB: 0--255, CMYK: 0--100\,\%, HSV: 0--360$^\circ$, 0--100\,\%).
    \item \textbf{Выбор из палитры и пипетки}: встроенная сетка дискретных цветовых образцов (swatches), диалоговое окно системной палитры \texttt{colorchooser}, поле ввода шестнадцатеричного кода (HEX) и генератор случайного цвета.
\end{enumerate}

\vspace{0.5cm}

\newpage
\noindent\textbf{4. Скриншот работающего приложения}

\begin{figure}[h!]
    \centering
    \includegraphics[width=0.88\textwidth]{app_screenshot.png}
\end{figure}

\vspace{0.5cm}
\noindent\textbf{5. Ключевые фрагменты программного кода}

\begin{lstlisting}[language=Python, caption={Математические функции взаимного преобразования}]
def rgb_to_cmyk(r, g, b):
    if r == 0 and g == 0 and b == 0:
        return 0.0, 0.0, 0.0, 1.0
    r_, g_, b_ = r / 255.0, g / 255.0, b / 255.0
    k = 1.0 - max(r_, g_, b_)
    c = (1.0 - r_ - k) / (1.0 - k)
    m = (1.0 - g_ - k) / (1.0 - k)
    y = (1.0 - b_ - k) / (1.0 - k)
    return c, m, y, k

def cmyk_to_rgb(c, m, y, k):
    r = 255.0 * (1.0 - c) * (1.0 - k)
    g = 255.0 * (1.0 - m) * (1.0 - k)
    b = 255.0 * (1.0 - y) * (1.0 - k)
    return r, g, b

def hsv_to_rgb(h, s, v):
    h = h % 360.0
    c = v * s
    x = c * (1 - abs((h / 60.0) % 2 - 1))
    m = v - c
    if h < 60:    rp, gp, bp = c, x, 0
    elif h < 120: rp, gp, bp = x, c, 0
    elif h < 180: rp, gp, bp = 0, c, x
    elif h < 240: rp, gp, bp = 0, x, c
    elif h < 300: rp, gp, bp = x, 0, c
    else:         rp, gp, bp = c, 0, x
    return (rp + m) * 255.0, (gp + m) * 255.0, (bp + m) * 255.0
\end{lstlisting}

\vspace{0.5cm}
\noindent\textbf{6. Вывод}

В ходе выполнения лабораторной работы были изучены аддитивные, субтрактивные и перцептивные принципы формирования цвета. Разработано графическое приложение, реализующее мгновенную синхронизацию между моделями CMYK, RGB и HSV тремя независимыми способами ввода. Устранены краевые эффекты нормализации угла $H$ и зацикливания обратного расчёта CMYK.

\end{document}